% HW1 brief: COMS E6998-019, Design of Production Agentic Systems.
% Build: cd briefs && tectonic -X compile -Z shell-escape -Z shell-escape-cwd=. hw1.tex
% minted needs pygmentize on PATH (e.g. uv tool install pygments).
\documentclass[9pt,letterpaper]{article}
\usepackage[extreme]{savetrees}
\usepackage{booktabs,tabularx,array,enumitem,xcolor}
\usepackage[cache=false]{minted}
\usepackage[breakable]{tcolorbox}
\linespread{1.15}
\tcbuselibrary{skins} % the preamble's `enhanced' key needs this library
\usepackage{xltabular} % rubric tables may break across pages
\usepackage[hidelinks]{hyperref}
\setlist{nosep,leftmargin=*}
\renewcommand\tabularxcolumn[1]{>{\raggedright\arraybackslash}p{#1}}
\setminted{fontsize=\small,breaklines,frame=single,framerule=0.3pt,framesep=2mm}
\usepackage{needspace}
% Listings never split across pages: each one, with its title, sits in an
% unbreakable box and moves to the next page whole if it does not fit.
\usepackage{etoolbox}
\BeforeBeginEnvironment{minted}{\par\smallskip\noindent\begin{minipage}{\linewidth}}
\AfterEndEnvironment{minted}{\end{minipage}\par\smallskip}
\newcommand\filelisting[3]{% title, lexer, file
  \par\smallskip\noindent\begin{minipage}{\linewidth}%
  \centerline{\small #1}\smallskip\inputminted{#2}{#3}\end{minipage}\par\smallskip}
\newtcolorbox{required}[1]{enhanced,breakable,colback=red!3,colframe=red!60!black,
  title={#1},fonttitle=\bfseries}
\title{HW1: A Bounded Agent for Bandit Alert Classificiation}
\author{COMS E6998 Section 019, Fall 2026 · Design of Production Agentic Systems\\
  Instructor: Rahul Krishna · TA: In Keun Kim}
\date{Released Friday, September 25, 2026 · Due Friday, October 16, 2026, before 11:59 p.m.}
\begin{document}
% savetrees' tight paragraphs set \looseness=-1 on every paragraph, which breaks
% ragged table cells early. Keep normal line breaking.
\makeatletter\markeverypar{\the\everypar}\makeatother
\maketitle

% Inline code: prints its argument verbatim in typewriter type (no # or % inside).
\newcommand\code[1]{\texttt{\detokenize{#1}}}
% File paths: typewriter type, breakable after / . _ (no spaces inside).
\newcommand\file{\nolinkurl}

\section{Overview}

You will build a vulnerability management agent. HW1 is limited to agentic
scanning of a Python repository and classification of the alerts you find. 

The HW handout fixes the interfaces you'll be designing to make them consistent and testable:
\begin{itemize}
\item \textbf{The command line}: how you invoke your agent from the terminal.
\item \textbf{The configuration}: the settings your agent reads from the config file.
\item \textbf{The output file}: where your agent writes the results of its analysis.
\item \textbf{The trace}: the detailed record of the agent's execution.
\item \textbf{The README}: the documentation that explains how to use your agent.
\end{itemize}

You design and build everything behind them. You get one point for 
each stage of the agent, so your grade shows which mechanisms exist and work.

Grades assess only the implementation of the required agent mechanisms, not
classification correctness or performance. Classification results are collected
for teaching purposes only. Submit the design report and implementation together;
the report and README document and provide evidence of the implementation.
There is no separate design checkpoint. You may consult the instructors as often
as you like, at any stage of the work.

\begin{center}
\begin{tabularx}{\linewidth}{@{}lX@{}}
\toprule
Released & Friday, September 25, 2026 \\
Due & Friday, October 16, 2026, before 11:59 p.m. \\
Weight & 20\% of the course grade, marked out of 20 (Section~\ref{sec:rubric}) \\
Format & Solo/Team assignment. One submission per student/team. \\
Late days & 6 per student per semester, at most 3 on one assignment; each extends the deadline by 24 hours. \\
\bottomrule
\end{tabularx}
\end{center}

\section{What you build}

Your agent takes one (pinned) Python repository and works in two stages.
\begin{enumerate}

  \item \textbf{Discovery.} Run a tool called bandit (Python security linter, its a CLI tool, see \url{https://bandit.readthedocs.io/en/latest/}).
Set it at \textbf{MEDIUM or higher severity} over the scan scope that the configuration sets for you (Section~\ref{sec:scope}).

\item \textbf{Classification.} Label every alert \textbf{1} (potential vulnerability)
  or \textbf{0} (not a vulnerability), with a text rationale block that includes line evidence.
\end{enumerate}

You then write a \code{report.json} file, with one row per alert, and record the agent's trace in a \code{trace.jsonl} file, the record of the run.

You write all of the code. The handout gives no starter code, stubs or tests.
Any language and framework is allowed, but Bandit needs Python 3.12. The target
code is read only: your agent reads it, and never installs or runs it.

We ship these files in this repository: \code{config.schema.json},
\code{report.schema.json} and \code{trace.schema.json} at the root; the three
required \code{[scan]} blocks as \file{briefs/scan/radicale.toml},
\file{briefs/scan/sglang.toml} and \file{briefs/scan/owasp-benchmark-python.toml};
and the example output \file{briefs/report.example.json}.

\section{Targets}\label{sec:targets}

Clone each target at its pinned commit. The root of that clone is
\code{--input}. Your agent must run on all three targets.

\begin{center}
\small
\newcommand\repo[2]{\begin{tabular}[t]{@{}l@{}}\footnotesize\url{#1}\\\ttfamily\footnotesize #2\end{tabular}}
\begin{tabularx}{\linewidth}{@{}lXlr@{}}
\toprule
Target & Repository and commit & Role & Alerts \\
\midrule
OWASP BenchmarkPython & \repo{https://github.com/OWASP-Benchmark/BenchmarkPython}{f1291485808b66e20ddb6b01b10dc71b3df8c8ba} & development & 18 alerts \\
Radicale & \repo{https://github.com/Kozea/Radicale}{eff8027f3dc4910be1659ee71f4b4a454ade5a8c} & held out & 10 alerts \\
SGLang & \repo{https://github.com/sgl-project/sglang}{bbe9c7eeb520b0a67e92d133dfc137a3688dc7f2} & held out & 21 alerts \\
\bottomrule
\end{tabularx}
\end{center}

\textbf{OWASP BenchmarkPython is the development set.} Its answer key is open to
you: the benchmark's own labels, in the file \code{expectedresults-0.1.csv} at
the repository root. Read it to measure your agent and to refine its
prompts, tools and stopping rules. \textbf{Your agent may not read it.} At run
time the key is withheld from the agent's workspace
(Section~\ref{sec:workspace}); an agent that reads it measures nothing.

\textbf{Radicale and SGLang are held out.} Their answer keys are hidden, and
staff evaluate classification results after submission for teaching purposes
only, with no effect on grades. They show whether your agent generalises or has
only been tuned to the benchmark.

\section{Required scan scope for each target}\label{sec:scope}

\begin{required}{Required scan scope for each target}
These three \code{[scan]} blocks are mandatory. Use them unchanged: graders run
every submission with exactly these blocks. \code{paths} and \code{exclude} are
globs relative to \code{--input}, and your Bandit wrapper expands them.

\filelisting{Radicale: 10 alerts}{toml}{scan/radicale.toml}
\filelisting{SGLang: 21 alerts}{toml}{scan/sglang.toml}
\filelisting{OWASP BenchmarkPython: 18 alerts (test cases 00001 to 00099)}{toml}{scan/owasp-benchmark-python.toml}
\end{required}

The counts come from Bandit 1.9.4 on CPython 3.12, measured on fresh clones at
the pinned commits. They count results of MEDIUM or higher severity, at any
confidence.

\subsection*{From a \code{[scan]} block to Bandit flags}

Your Bandit wrapper turns each key into a command-line flag. Run Bandit from
the root of the target checkout (the \code{--input} directory).

\begin{center}
\begin{tabularx}{\linewidth}{@{}lX@{}}
\toprule
\code{[scan]} key & Bandit command line \\
\midrule
\code{paths} & Positional arguments, after glob expansion; \code{-r} recurses into directories. \\
\code{exclude} & \code{-x} with a comma-separated list; omit the flag when the list is empty. \\
\code{min_severity = "MEDIUM"} & \code{--severity-level medium} (the same as \code{-ll}). \\
(any confidence) & No \code{--confidence-level} flag; the default reports every confidence. \\
\code{bandit_version} & Install \code{bandit==1.9.4}; \code{bandit --version} must report it. \\
(output) & \code{-f json -o <file>}, which your wrapper parses. \code{-q} suppresses the banner. \\
\bottomrule
\end{tabularx}
\end{center}

The three required blocks translate to these commands. Each gives the alert
count stated above. Bandit exits with status 1 when it reports results, so do
not treat that exit status as a failure.

\begin{minted}{bash}
# Radicale (from targets/radicale)
bandit -q -r --severity-level medium -x radicale/tests -f json -o bandit.json radicale

# SGLang (from targets/sglang)
bandit -q -r --severity-level medium -f json -o bandit.json \
  python/sglang/multimodal_gen/runtime scripts/playground/replay_request_dump.py

# OWASP BenchmarkPython (from targets/owasp-benchmark-python); the shell expands the glob
bandit -q -r --severity-level medium -f json -o bandit.json testcode/BenchmarkTest000*.py
\end{minted}

Bandit adds a leading \code{./} to the paths of files named on the command
line, as in the OWASP command. Strip that prefix before you write
\code{report.json} (Section~\ref{sec:output}).

\section{Interface}\label{sec:interface}

The agent is a command-line program:
\begin{minted}{bash}
./your-agent --input <target checkout root> --output <out dir> --config config.toml
\end{minted}
It writes two files into \code{<out dir>}: \code{report.json}
(Section~\ref{sec:output}) and \code{trace.jsonl} (Section~\ref{sec:trace}). You choose the name
\code{your-agent}, and your README states it.

Your agent must read the configuration file named by \code{--config}. It is a
TOML file, and once parsed it must validate against \code{config.schema.json}.
You ship two configs, a full budget and a low budget (Section~\ref{sec:readme}),
and we run your agent with them.

\begin{center}
\small
\begin{tabularx}{\linewidth}{@{}lX@{}}
\toprule
Key & Rule \\
\midrule
\code{[scan]} & Exactly one of the three required blocks (Section~\ref{sec:scope}), unchanged. \\
\code{model.provider}, \code{model.name} & Required strings: your provider and the exact model name. \\
\code{budget.max_tokens} & Required integer $\geq 0$: the token budget for the whole run, across all alerts. \\
\code{budget.output_allowance} & Required integer $\geq 1$: the maximum output tokens for each model call, sent to the provider. \\
\code{budget.max_steps} & Required integer $\geq 1$: a failsafe, separate from the token budget. \\
(other keys) & Allowed, for your own settings. Never put secrets in the file; read them from environment variables. \\
\bottomrule
\end{tabularx}
\end{center}

A valid config for Radicale:
% begin config example (the staff release checks validate this block against config.schema.json)
\begin{minted}{toml}
[scan]
paths = ["radicale"]
exclude = ["radicale/tests"]
min_severity = "MEDIUM"
bandit_version = "1.9.4"

[model]
provider = "anthropic"
name = "claude-sonnet-5"

[budget]
max_tokens = 200000
output_allowance = 1024
max_steps = 200
\end{minted}
% end config example

The token budget covers the whole run, not each alert. The budget charges
\code{input + output} as the provider reported them, summed over every
\code{model_call} (Section~\ref{sec:trace}). \code{cache_read} and
\code{cache_write} are recorded when the provider reports them, and are never
charged. \code{budget.used_tokens} in \code{report.json} is exactly that sum.
When the budget runs out, every alert that is not yet classified still gets a
row in \code{report.json}, with \code{status} set to \code{budget_exhausted}.

\section{Output}\label{sec:output}

\code{report.schema.json} (JSON Schema) is normative: your \code{report.json}
must validate against it. The example below is a truncated illustration: it shows 3 of the 18 rows of an
OWASP BenchmarkPython cutoff run, and its \code{alert_count} equals its 3 rows.
It shows the format only; the rationales are placeholders.
\filelisting{briefs/report.example.json (truncated illustration: 3 of the 18 OWASP rows)}{json}{report.example.json}

Rules:
\begin{itemize}
\item There is exactly one row for each in-scope Bandit result. Rows are never
  dropped or merged, so \code{len(alerts)} equals \code{scan.alert_count}.
\item \code{path} and every evidence \code{path} are repository-relative POSIX
  paths with no leading \code{./} (Bandit prints \code{./} for explicitly listed
  files; strip it). Rows are ordered by plain string comparison of \code{path},
  then \code{line} as an integer, then \code{test_id}.
\item \code{status} is one of \code{classified}, \code{budget_exhausted} or \code{error}.
\item \code{label} is \code{1} or \code{0} when \code{status} is \code{classified}.
  Otherwise it is \code{null}.
\item \code{budget.used_tokens} is exactly the budget charge
  (Section~\ref{sec:interface}): \code{input + output} as the provider reported
  them, summed over every \code{model_call}. \code{cache_read} and
  \code{cache_write} are never charged.
\end{itemize}

\section{Workspace and withholding}\label{sec:workspace}

You build the workspace and the withholding yourselves; there is no deny list
in the configuration. The reference pattern is cyberbird's
\file{cyberbird/reactive/workspace.py}, in particular \code{resolve},
\code{AgentWorkspace} and \code{WITHHELD_GLOBS}:
{\raggedright\url{https://github.com/comse6998-019/cyberbird/blob/cf7e96c8bd9e74470b99b7b754bb72392427cfc6/cyberbird/reactive/workspace.py}\par}

Discovery (the Bandit run) scans the whole required \code{[scan]} scope. The
withholding rules apply to what the model can read through the classificiation tools:
read, search, and any other tool the model can call.

Requirements:
\begin{enumerate}
\item The agent works in a disposable copy of \code{--input}. The tools the model
  can call reach only that copy.
\item Every tool path goes through one path check. The check refuses \code{..}
  paths, absolute paths, and symlinks that resolve outside the copy.
\item \code{.git/} is withheld on every target. The history of a full clone can
  hold commits made after the pinned one, including fixes.
\item On OWASP BenchmarkPython, the answer key \code{expectedresults-0.1.csv}
  (in general, \code{expectedresults*.csv}) is withheld from the agent. So is
  every \code{BenchmarkTest*} file except the case under investigation, because
  each case has a safe twin. You may still read the key yourselves
  (Section~\ref{sec:targets}); only the agent's view is restricted.
\end{enumerate}

\section{Trace and trajectory}\label{sec:trace}

The trace is the agent's trajectory: the complete, ordered record of what one
run did and who decided each step. Graders read it to check B3 to B8, so its
format is fixed. You build it yourselves. The reference implementation is
cyberbird's \file{cyberbird/plan_and_validation/trace.py} (\code{Trace},
\code{EventKind}, \code{TerminalStatus}, \code{Usage}) and
\file{cyberbird/plan_and_validation/trajectory.py} (trace to Mermaid sequence
diagram):
{\raggedright\url{https://github.com/comse6998-019/cyberbird/blob/cf7e96c8bd9e74470b99b7b754bb72392427cfc6/cyberbird/plan_and_validation/trace.py}\\
\url{https://github.com/comse6998-019/cyberbird/blob/cf7e96c8bd9e74470b99b7b754bb72392427cfc6/cyberbird/plan_and_validation/trajectory.py}\par}

Format:
\begin{itemize}
\item \code{<out dir>/trace.jsonl} holds one JSON object per line. It is
  append-only and flushed after every event, so that a crashed run leaves the
  trace that explains the crash.
\item Every event has \code{step} (an integer from 1, increasing by 1),
  \code{ts} (Unix time, float), \code{run_id}, \code{alert_id} (the alert the
  event belongs to, or \code{null} for run-level events) and \code{kind}.
\item \code{kind} is one of the closed set below. Any event may carry other
  fields as well.
\end{itemize}

\begin{center}
\small
\begin{tabularx}{\linewidth}{@{}lp{0.25\linewidth}X@{}}
\toprule
\code{kind} & Written by & Required fields \\
\midrule
\code{model_call} & the node that called the model & \code{role}, \code{model},
  \code{usage} with \code{input} and \code{output} (integers, as the provider
  reported them; also \code{cache_read} and \code{cache_write} if the provider
  reports them; no summed total) \\
\code{tool_request} & the dispatcher & \code{tool}, \code{args} \\
\code{tool_result} & the dispatcher & \code{tool}, \code{ok} (boolean), and
  \code{error} when \code{ok} is false \\
\code{routing} & the edge function that chose the next node & \code{decision},
  \code{by} (\code{model} or \code{runtime}) \\
\code{state_change} & any node, for updates not covered above & \code{node};
  when an alert's row is settled (\code{classified}, \code{error}, or
  \code{budget_exhausted}, including rows settled at a crash or at the cutoff),
  also \code{alert_status} and \code{label} \\
\code{terminal} & the runtime, once, last & \code{status} (one of your terminal
  statuses, A6), \code{usage_total} (an object with the same counters as \code{usage}) \\
\bottomrule
\end{tabularx}
\end{center}

Invariants (graders check these on every trace):
\begin{enumerate}
\item The file parses line by line, and every event has the required fields for
  its \code{kind}.
\item \code{step} starts at 1 and increases by exactly 1.
\item There is exactly one \code{terminal} event, and it is the last line, also
  for a run that crashed.
\item Each counter in \code{usage_total} equals the sum of that counter over all
  \code{model_call} events. \code{budget.used_tokens} in \code{report.json}
  equals exactly the sum of \code{input + output} over those events;
  \code{cache_read} and \code{cache_write} are recorded but never charged.
\item For every row in \code{report.json}, including \code{budget_exhausted} and
  crash-time \code{error} rows, there is exactly one\linebreak{} \code{state_change} event
  with that row's \code{alert_id} whose \code{alert_status} and \code{label}
  equal the row's \code{status} and \code{label}.
\item Every \code{tool_request} is immediately followed by its \code{tool_result}
  (the dispatcher handles one request at a time), before the next
  \code{model_call}.
\end{enumerate}

\code{trace.schema.json}, at the root of this repository, is normative for each
line: every line of your \code{trace.jsonl} must validate against it. It encodes
the kind table above; the invariants are rules across lines, which a schema
cannot express, so graders check them separately. \code{briefs/trace.example.jsonl}
is the trajectory behind \code{briefs/report.example.json}: it validates line by
line and meets all six invariants.


\section{Rubric}\label{sec:rubric}

The assignment is marked out of 20, one point per item. These items assess the
implementation of the required mechanisms and the documentation and evidence
that demonstrate them, not the correctness or performance of classifications.

{\small
\begin{xltabular}{\linewidth}{@{}lX@{}}
\toprule
\multicolumn{2}{@{}l}{\textbf{Part A: design report (7)}} \\
\midrule
A1 & \textbf{Behaviours.} What the agent does for each alert, and what it must
  never do: run target code, write to the target, invent observations. \\
A2 & \textbf{State machine.} A diagram of the states and the transitions between them. \\
A3 & \textbf{Nodes and edges.} Each node, each edge marked as conditional or
  fixed, and who decides each transition (the model or the runtime). \\
A4 & \textbf{Agent state.} Each field with its reducer, its lifetime, who writes
  it, and whether the model sees it. \\
A5 & \textbf{Tools.} Contracts for at least 3 tools, with arguments, limits and refusal cases. \\
A6 & \textbf{Start, stop and termination.} A closed, mutually exclusive set of
  terminal statuses and what triggers each one. This includes the budget
  policy: how the charge above is applied, when it is checked, how the output
  allowance is sent, and what happens at the cutoff. \\
A7 & \textbf{Workspace and withholding.} Which files the agent must not see for
  each target and why, how they are withheld, and how they are restored. \\
\bottomrule
\end{xltabular}

\begin{xltabular}{\linewidth}{@{}lX@{}}
\toprule
\multicolumn{2}{@{}l}{\textbf{Part B: implementation (10)}} \\
\midrule
B1 & \textbf{Bandit tool.} Reads \code{[scan]}, uses MEDIUM+ severity, uses
  Bandit 1.9.4, and parses the JSON output. \\
B2 & \textbf{Workspace, read and search tools.} Meets every requirement in
  Section~\ref{sec:workspace}. \\
B3 & \textbf{Dispatcher.} The runtime validates and runs every request. Unknown
  tools and bad arguments come back as refusal observations, not crashes. Only
  the runtime writes observations. \\
B4 & \textbf{Graph.} Implemented as the design in Part A describes. \\
B5 & \textbf{Tracer and trajectory.} Writes \code{trace.jsonl} in the
  Section~\ref{sec:trace} format with every invariant holding: one event
  per step, flushed immediately, routing events that record who decided, and
  exactly one terminal event, even after a crash. \\
B6 & \textbf{Budget capture.} Charges the usage the provider reports, not an estimate. \\
B7 & \textbf{Budget cutoff.} Checked before each model call. Alerts that are not
  yet classified get \code{budget_exhausted}. \\
B8 & \textbf{Failsafes.} \code{max_steps} and a no-progress detector, each ending
  the run with its own terminal status. \\
B9 & \textbf{\code{report.json}.} Validates against \code{report.schema.json}
  with one row per alert, and is still written when the run crashes. \\
B10 & \textbf{All three targets.} Runs end to end with the staff \code{[scan]} blocks. \\
\bottomrule
\end{xltabular}

\begin{xltabular}{\linewidth}{@{}lX@{}}
\toprule
\multicolumn{2}{@{}l}{\textbf{Part C: evidence and README (3)}} \\
\midrule
C1 & \textbf{README.} Following it from a fresh clone
  (Section~\ref{sec:readme}), we can run your agent end to end on each target
  with both of your configs, and the low-budget run shows the cutoff. \\
C2 & \textbf{Results.} For each target, the counts of label 1, label 0 and
  \code{null}. On OWASP, precision and recall against the benchmark labels, and
  how you used them to refine the agent (metrics are reported for teaching
  purposes only; their values are not graded). One
  sequence diagram of one alert's trajectory, drawn from one of your traces. \\
C3 & \textbf{Limitations and alternatives.} A comparison with at least two other
  architectures (fixed workflow, plan-and-execute, supervisor/worker) on cost,
  latency, coordination and failure behaviour. Analyse them; do not build them. \\
\bottomrule
\end{xltabular}
}

Grading notes:
\begin{itemize}
\item Grades assess implementation only; the design report, README and traces
  provide supporting evidence. Classification correctness and performance
  (including accuracy, precision and recall) are evaluated only for teaching
  purposes and do not affect grades.
\item A missing mechanism cannot earn its point.
\item We check B7 with your \file{configs/low.toml}.
\end{itemize}

\section{README}\label{sec:readme}

Your \code{README.md} tells us how to run your agent. We follow it from a fresh
clone of your repository, point your agent at each target checkout, and run it
end to end. How you install and invoke the agent is up to you; the README only
has to let us do it without asking you.

It must give:
\begin{enumerate}
\item \textbf{Setup} from a fresh clone: every command to install your
  dependencies, Python 3.12 and Bandit 1.9.4, and the environment variables your
  agent reads (names only, never values).
\item \textbf{The run command}, in the form of Section~\ref{sec:interface}, with
  \code{--input}, \code{--output} and \code{--config}.
\item \textbf{Two config files}, committed in your repository, that validate
  against \code{config.schema.json}:
  \begin{itemize}
  \item \file{configs/full.toml}: the full budget you consider right for a
    complete run;
  \item \file{configs/low.toml}: a \code{max_tokens} low enough that the budget
    runs out partway through the Radicale run, so at least one row of
    \code{report.json} has \code{status} \code{budget_exhausted}.
  \end{itemize}
  Put the Radicale \code{[scan]} block in both. For each target we replace the
  \code{[scan]} block with that target's required block
  (Section~\ref{sec:scope}) and keep everything else.
\end{enumerate}

We run your agent six times: on each of the three targets, with each of the two
configs. Each run must write \code{report.json} and \code{trace.jsonl}
(Sections~\ref{sec:output} and~\ref{sec:trace}). Anything else in the README is
up to you.

\section{Deliverables}
\begin{itemize}
\item Source code, and a lockfile or equivalent that pins dependencies.
\item \code{README.md} as Section~\ref{sec:readme} describes, with
  \file{configs/full.toml} and \file{configs/low.toml}.
\item The design report (Part A plus C2 and C3), as \code{REPORT.md} or a PDF.
\item The unedited \code{report.json} and \code{trace.jsonl} of every run
  your report cites, with the config used. No API keys in the repository.
\item A contribution statement that says who did what.
\end{itemize}

\section{Submissions}

We will create the HW1 assignment on CourseWorks (Canvas). Submit there, and
follow the instructions on that assignment page. Each team makes one
submission: any one member submits it on behalf of the whole team.

\section{Academic integrity and AI tools}

AI assistants may be used to learn, debug your own code, review a design, or
provide scoped autocomplete. Do not hand the assignment to an agent and submit
the result. You remain responsible for every line you submit. Every team member
must be able to explain the entire submission; the individual exams test this.

You may not use agents to generate or modify the target code. The guideline on
LLM use is strictly enforced: use LLMs to learn, not to do the assignment for
you.

\end{document}
